\section{Related Work}
\subsection{Retrieval Augmented Generation}

Large Language Models (LLMs) demonstrate strong capabilities in generating text, but there are also obstacles such as outdated information and the potential to hallucinate~\cite{tonmoy2024comprehensive, wang2023survey, huang2023survey}. To address these issues, RAG retrieves external knowledge to generate responses with improved accuracy and factuality~\cite{gao2023retrieval, yu2024evaluation, huang2024survey}. Integrating external knowledge is especially crucial in fields like legal, medical and finance, where precision and reliability are essential~\cite{louis2024interpretable, xiong2024benchmarking, zhang2023enhancing}.

RAG systems have shown impressive performance across a range of tasks, including open-domain question answering~\cite{mao2020generation, he2024g, lewis2020retrieval}, code generation~\cite{parvez2021retrieval, zhou2022docprompting, tan2024prompt} and dialogue~\cite{shuster2021retrieval, komeili2021internet, thulke2021efficient}. Additionally, real world products like Bing Search\footnote{\url{https://www.bing.com/chat}} and Langchain~\cite{langchain} have integrated applications based on RAG.

\subsection{Evaluation of RAG}
\label{sec:rw_evaluation_of_rag}
Existing evaluation practices for RAG systems can be categorized into two main approaches: evaluating essential capabilities of generators only and assessing end-to-end performance of RAG systems.

Within the two components of a RAG system, the retriever has been well studied in recent years, thus a line of recent work focused on evaluating essential generator capabilities. 
RGB~\cite{chen2024benchmarking} evaluated 4 fundamental abilities required for generators including Noise Robustness, Negative Rejection, Information Integration and Counterfactual Robustness by manually constructed test sets. 
RECALL~\cite{liu2023recall} introduced manually edited counterfactual contexts into QA and text generation datasets to evaluate the counterfactual robustness of LLMs. 
NoMIRACL~\cite{thakur2023nomiracl} evaluated LLMs' robustness against first-stage retrieval errors of RAG systems with manually judged relevant and non-relevant datasets. 
Wu et al.~\cite{wu2024faithful} quantified the tug-of-war between LLMs’ faithfulness and internal prior by introducing varying levels of perturbations on the provided contexts.
FaaF~\cite{katranidis2024faaf} introduced a fine-grained fact verification formulation to improve previous prompting-based approaches in evaluating factuality of generators.
However, we argue that above generator-only evaluation approaches with manually constructed datasets cannot serve as a general RAG evaluation framework to reveal the entanglement of between generation results and different retrieval behaviors, as shown in the analysis of Sec.~\ref{sec:analysis}.

Another line of work focused on assessing end-to-end quality scores of RAG systems. TruLens~\cite{TruLens} introduced the concept of RAG Triad, which decompose the quality scores into three aspects: context relevance, groundedness and answer relevance, then predicted the score by prompting LLMs or using NLI models. RAGAS~\cite{es2023ragas} and ARES~\cite{saadfalcon2023ares} followed the RAG Triad concept and improved the score prediction approaches on different datasets. CRUD-RAG~\cite{lyu2024crud} refered to the CRUD (Create, Read, Update and Delete) actions between users and knowledge bases to develop corresponding datasets and evaluation metrics for RAG systems. We compare the above four evaluation frameworks with \benchname\space in the meta evaluation of Sec.~\ref{sec:meta_evaluation}.

Besides, the following work also provided good insight or high quality datasets for end-to-end RAG evaluation. Liu et al.~\cite{liu2023evaluating} conducted human evaluation to audit four popular generative search engines in terms of fluency, perceived utility, and verifiability. MEDRAG~\cite{xiong2024benchmarking} constructed a medical RAG benchmark from medical QA datasets and evaluated medical RAG systems with QA accuracy. MultiHop-RAG~\cite{tang2024multihop} generated multi-hop queries from news articles and evaluated RAG systems with QA accuracy. CDQA~\cite{xu2024let} proposed a novel approach to generate dynamic QA questions which requires latest information to answer. However, the evaluation metrics used in the work mentioned above rely either on human evaluation or simple textual accuracy, making them incapable of complex RAG scenarios that require long answer evaluation. Therefore, we do not include them in the meta evaluation.
